%! The Rational Parent Function \& Transformations — Unit 4, Lesson 4.4 — Warm-Up (Answer Key)
\documentclass[10pt]{article}
\usepackage{algebra2-article}
\usepackage{algebra2-key}   % pulls in -boxes; do NOT also load -boxes
\newcommand{\TallMath}[1]{$\displaystyle #1\rule[-1.4em]{0pt}{3.2em}$}

\begin{document}

\pageheader{Unit 4, Lesson 4.4}{Warm-Up --- Answer Key}
\namedateperiod

\vspace{0.10in}

\begin{notesbox}{Getting Ready: Skills We'll Use Today}
\small Three skills, one per leg of today's lesson. Work quickly.

\vspace{4pt}
\textbf{1. The reciprocal, up close and far out} (Lesson 4.0). Complete the table for
$y=\dfrac{1}{x}$.
\begin{center}
\renewcommand{\arraystretch}{1.35}
\begin{tabular}{|c|c|c|c|c|c|c|c|}
\hline
$x$ & $-4$ & $-1$ & $-\frac12$ & $0$ & $\frac12$ & $1$ & $4$ \\ \hline
$y=\frac1x$ & \ans{$-\frac14$} & \ans{$-1$} & \ans{$-2$} & \ans{undef.} & \ans{$2$} & \ans{$1$} & \ans{$\frac14$} \\ \hline
\end{tabular}
\end{center}
\begin{enumerate}[label=(\alph*), leftmargin=2.1em, itemsep=3pt]
  \item Domain: \ans{all reals except $x=0$} \quad Range: \ans{all reals except $y=0$}
  \item As $x$ gets close to $0$ the outputs \ans{blow up (grow without bound)}, so the graph has a \emph{vertical
        asymptote} with equation \ans{$x=0$}.
  \item As $x$ runs far out to the right or left the outputs \ans{settle toward $0$}, so the graph has a
        \emph{horizontal asymptote} with equation \ans{$y=0$}.
\end{enumerate}

\vspace{4pt}
\textbf{2. Transformations you already own} (Units 1--2). Name the turning point of each.
\begin{enumerate}[label=(\alph*), leftmargin=2.1em, itemsep=3pt]
  \item $y=|x-3|+1$: \; vertex $($\ans{$3$}$,\,$\ans{$1$}$)$ --- the parent moved
        \ans{right} $3$ and \ans{up} $1$.
  \item $y=(x+2)^2-5$: \; vertex $($\ans{$-2$}$,\,$\ans{$-5$}$)$ --- the parent moved
        \ans{left} $2$ and \ans{down} $5$.
  \item Fill in the rule you have used since Unit 1: a number attached \emph{outside} the function
        moves the graph \ans{vertically} and does exactly what it says; a number attached
        \emph{inside} moves the graph \ans{horizontally} and does the \ans{opposite} of what it says.
\end{enumerate}

\vspace{4pt}
\textbf{3. One function, two costumes} (Lessons 4.3 and 4.0).
\begin{enumerate}[label=(\alph*), leftmargin=2.1em, itemsep=3pt]
  \item Combine into a single fraction: \; $2+\dfrac{3}{x-1}=\dfrac{2(\text{\ans{$x-1$}})+3}{x-1}
        =$ \ans{$\frac{2x+1}{x-1}$}
  \item For that single fraction, set the denominator to zero: vertical asymptote
        $x=$ \ans{$1$}.
  \item Its numerator and denominator have the \emph{same} degree, so by the degree comparison the
        horizontal asymptote is $y=$ \ans{$2$} (the ratio of the leading coefficients).
\end{enumerate}

\end{notesbox}

\begin{teachernote}
Each item is one leg of today's lesson. \textbf{Item 1} is the parent function itself, met numerically
before it is met as a picture --- and it is the item to debrief slowly, because two of today's answers
are permanent: this parent has \emph{no} $x$-intercept and \emph{no} $y$-intercept, which is unlike every
parent students have graphed so far. If someone writes $y=0$ in the $x=0$ column, that is the whole
lesson's first correction.

\vspace{3pt}
\textbf{Item 2} is the transformation rule from Units 1 and 2, stated in the general form students will
apply to a brand-new parent today. Insist on 2(c) in words. Today the ``inside is backwards'' rule is
worth more than it has ever been, because on this parent an inside shift moves a \emph{vertical
asymptote} --- and a student who thinks $x-3$ shifts left will put the wall on the wrong side of the
$y$-axis and never recover.

\vspace{3pt}
\textbf{Item 3} is the payoff to watch, and it is the one to say the least about. Students combine
$2+\frac{3}{x-1}$ into $\frac{2x+1}{x-1}$, then find its asymptotes the Lesson 4.0 way --- and the
answers come out $x=1$ and $y=2$, which are exactly the two numbers they started with. Do not name the
coincidence now. It is Section 6 of the notes, where the class discovers that a rational function whose
top and bottom are both linear \emph{is} a shifted, stretched $\frac1x$ in disguise --- and that the
$k$ they add outside \emph{is} the horizontal asymptote.
\end{teachernote}

\end{document}
